\section{Limitations and Claim Boundaries}\label{sec:limitations}

The dataset is small, single-application, and structurally homogeneous.  Its 33
directories are not an independent random sample, and the 27 complete traces
have identical event totals.  Percentages are descriptive only.  Six
source-only directories have unknown outcomes, introducing survivor bias if
analysis is restricted to complete sessions.

No controlled baseline, adversarial campaign, fault injection, concurrency
study, performance benchmark, or usability study was conducted.  The bounded
prompt-content field was present in 0/27 complete valid traces; six source-only
directories were unassessable.  Thus prompt-level qualitative analysis and
injection measurement are impossible.  We do not infer attack
success from policy violations: those violations report unknown rules, not
successful attacks.

Offline verification checks bytecode, trace, security-report, and
trace-event-ledger digest consistency against an unsigned local bundle.  It does not
check source integrity, authenticate authorship, establish trusted time, prevent
replacement of the entire bundle, or prove that the trace is complete.  SHA-256
is used for semantic digests; no post-quantum security property is evaluated
or claimed.

Passports, \code{ans_name}, jurisdictions, residency, ATrust identities,
credential contracts, handshakes, ledgers, MCP bridges, and A2A bridges are
partly or wholly declarative.  The current system does not resolve Sovereign
DNS/ANS names or DIDs, authenticate real agents, verify credentials, execute
cryptographic handshakes, provide blockchain consensus or immutable storage,
or run live MCP/A2A interoperability.

Application-level deny flags and the provider registry do not establish
production isolation.  We did not audit dependencies, prove memory safety of
all unsafe/transitive code, inspect operating-system calls, or test container
escapes.  We make no guarantee against all prompt injection, data leakage, or
malicious providers.

Governance mappings are not compliance certification.  Neither the language
nor this study certifies conformity with a law, regulator, NIST AI RMF, or any
other framework \citep{tabassi2023airmf}.  The empirical contradiction between
the top-level and detailed statuses further prevents a positive assurance
conclusion.

\begin{table}[t]
\centering
\caption{Limitations and future assurance requirements. These are excluded
from current claims unless the required mechanism and evaluation are added.}
\label{tab:limitations-assurance}
\begin{tabularx}{\linewidth}{@{}XXX@{}}
\toprule
Limitation & Current status & Future assurance requirement \\
\midrule
Live provider execution & Not implemented & sandboxed adapter plus host telemetry \\
Real identity authentication & Not implemented & DID/VC or equivalent resolver and issuer policy \\
Source integrity & Outside current bundle verification & source digest and source-to-bytecode binding \\
Producer authentication & Not implemented & signed bundles and key governance \\
Replacement detection & Not implemented & append-only log or transparency service \\
Large workload diversity & Not present & controlled workload matrix and pre-registered metrics \\
\bottomrule
\end{tabularx}

\end{table}

\Cref{tab:limitations-assurance} converts limitations into concrete future
requirements: source integrity needs source-to-bytecode binding, producer
authentication needs signatures and key governance, and replacement detection
needs an append-only transparency mechanism.

Future evaluations should publish controlled workloads, independent
verifiers, OS-level telemetry, and pre-registered outcome definitions.
